% -----------------------------------------------------------
\section{BOM Extended Scripture Studies}
% -----------------------------------------------------------
	\label{EXSS-BOM}
	
	% ----------------------
	\subsection{Spiritual manifestations and their contexts in the Book of Mormon}
	% ----------------------
		\label{EXSS-BOM-SPIRIT-CONTEXTS}
		
		This study was prompted by reading in the book of Ether, where Jared's brother sees Christ's spiritual body at the top of a mountain. It got me thinking about other great spiritual manifestations in the Book of Mormon, and the contexts in which they occurred.
		
		As I've been going through these, I've also realized that I need to divide the spiritual manifestations into larger and smaller ones. The larger ones are actual appearances or visions, and the smaller ones are just smaller ones.
		
		% -----
		\subsubsection{Larger manifestations}
		% -----
			
			\scriptsize

			\begin{longtable}{p{1cm}|p{1cm}|p{0.75cm}|p{1.5cm}|p{4cm}|p{4cm}|p{3.5cm}}
				
				% HEADER
				\textbf{\textsc{Who}} & \textbf{\textsc{Ref.}} & \textbf{\textsc{Prays?}} & \textbf{\textsc{For?}} & \textbf{\textsc{Environmental Context}} & \textbf{\textsc{Personal Context}} & \textbf{\textsc{What Happens?}} \\ \hline
				
				% Lehi and the rock
				Lehi &
				1 Nephi 1:5--6 &
				Yes &
				``in behalf of his people'' &
				\textbf{\tr{Outside, by a rock, alone.}} He seems to be outside of the city, and he prayed ``as he went forth.'' He seems to be alone. He's in a place with a rock large enough for a ``pillar of fire'' to rest on, so that's another reason to think he's outside of the city. After the manifestation, he ``returned to his own house at Jerusalem,'' and so couldn't have been at home, and the wording of that sentence also makes me think he was outside the city. &
				It's possible he listened to the prophets who were preaching and saying that Jerusalem would be destroyed unless the people repented. He's thinking about the people in the city and prays ``with all his heart, in behalf of his people.'' &
				He sees and hears ``much,'' but seems to stay mentally in the place he's at, by the rock. \\ \hline
				
				% Lehi on his bed
				Lehi &
				1 Nephi 1:7--15 & 
				No & 
				--- & 
				\textbf{\tr{At home in his room, alone.}} He's in his own house, on his own bed; this is after he returns to his place in Jerusalem. He's alone though. & 
				He's just had the prior vision and is now ``overcome with the Spirit,'' so he's got a lot going on in his head. He's also been overcome by the things he's seen, which before caused him to ``quake and tremble exceedingly.'' & 
				He gets ``carried away'' in this vision, and seems to see God and angels and others. So this one seems more out-of-body than his first one, at least. He ends up less upset than after this first vision, especially when he reads the book in verse 12 and is ``filled with the Spirit of the Lord.'' His soul rejoices in 15, etc. This vision also prompts him to begin to go out among the people of Jerusalem to preach repentance to them. \\ \hline
				
				% Summary
				 & % Who
				 & % Scripture reference
				 & % Prays?
				 & % For?
				 & % Environmental context
				 & % Personal context
				 \\ \hline % What happens?
				 
				% *LARGE
		
			\end{longtable}
		
			\normalsize
		
			\clearpage
			
		% -----
		\subsubsection{Smaller manifestations}
		% -----
		
			\scriptsize
		
			\begin{longtable}{p{1cm}|p{1cm}|p{0.75cm}|p{1.5cm}|p{4cm}|p{4cm}|p{3.5cm}}
				
				% HEADER
				\textbf{\textsc{Who}} & \textbf{\textsc{Ref.}} & \textbf{\textsc{Prays?}} & \textbf{\textsc{For?}} & \textbf{\textsc{Environmental Context}} & \textbf{\textsc{Personal Context}} & \textbf{\textsc{What Happens?}} \\ \hline
				
				% Lehi hearing he needs to go to the wilderness
				Lehi & % Who
				1 Nephi 2:1--2 & % Scripture reference
				No & % Prays?
				--- & % For?
				\textbf{\tr{In the city, at home, in a dream while asleep.}} He seems to be at home, and it's possible this one is at night, so maybe his wife and family are around, depending on their sleeping arrangements. So this is in the city, at home. & % Environmental context
				Lehi has been preaching, and apparently some threats have been made on his life. It's possible this has been on his mind a lot. & % Personal context
				Receives instruction in a dream. Note that it's possible verses 1 and 2 report different dreams. \\ \hline % What happens?
				
				% Nephi's heart is softened
				Nephi & % Who
				1 Nephi 2:16 & % Scripture reference
				Yes & % Prays?
				To know the mysteries of God & % For?
				\textbf{\tr{In the wilderness.}} All we know here is that Nephi is out in the wilderness with his parents and brothers. & % Environmental context
				Nephi had ``great desires to know of the mysteries of God,'' and so he felt a lack of knowledge and wanted to fill that lack. & % Personal context
				It isn't clear whether Nephi actually comes to know any ``mysteries'' in this manifestation; his heart gets softened so that he believes Lehi's words. In telling Sam about his experience, though, in 1 Nephi 2:17, he does say that the Lord manifested things to him ``by his Holy Spirit.'' \\ \hline % What happens?
				
				% Nephi prays for his brothers
				Nephi & % Who
				1 Nephi 2:18--24 & % Scripture reference
				Yes & % Prays?
				His brothers Laman and Lemuel, who won't listen to him, or to his father. & % For?
				\textbf{\tr{Away from camp, in the wilderness.}} Nephi here seems to be some distance away from the camp, because in 3:1, he says that he ``returned from speaking with the Lord to the tent of my father.'' So he was away off by himself. It could have been an overnight thing, because when he gets back, Lehi tells him that he's had an important dream. & % Environmental context
				He's upset that Laman and Lemuel wouldn't listen to him, and probably upset that they wouldn't listen to Lehi either. He says he is ``grieved because of the hardness of their hearts.'' & % Personal context
				The Lord gives him instruction and promises about what is going to happen with his brothers. \\ \hline% What happens?
				
				% Lehi dreams that his sons need to return to Jerusalem
				Lehi & % Who
				1 Nephi 3:2 & % Scripture reference
				No & % Prays?
				--- & % For?
				\textbf{\tr{In the wilderness, at camp.}} In the wilderness, while Nephi was away; we don't know whether it was from sleeping during the night, or if this was a Lehi nap that brought on the dream. & % Environmental context
				We don't really know, but perhaps he was thinking about the records and his ancestry. & % Personal context
				He shares the news with Nephi, Laman and Lemuel don't like it, etc. \\ \hline % What happens?
				
				% Angel appears in order to stop brothers fighting
				Nephi, Laman, Lemuel, Sam & % Who
				1 Nephi 3:29--30 & % Scripture reference
				No & % Prays?
				--- & % For?
				\textbf{\tr{In a cave outside of Jerusalem.}} & % Environmental context
				Laman and Lemuel are beating Nephi and Sam up. & % Personal context
				They're told to try again, though Laman and Lemuel start complaining immediately. Laman and Lemuel saw it too, though, as Nephi affirms in 4:3. \\ \hline % What happens?
				
				% Summary
				 & % Who
				 & % Scripture reference
				 & % Prays?
				 & % For?
				 & % Environmental context
				 & % Personal context
				 \\ \hline % What happens?
				 
				% *SMALL
				
			\end{longtable}
		
			\normalsize
		
			\clearpage
	
	% ----------------------
	\subsection{``Discernible'' in Alma 32:35}
	% ----------------------
		\label{EXSS-BOM-ALMA-32-DISCERNIBLE}
	
	This is the only use of ``discernible'' in all the scriptures. Here is the entry from the 1828 dictionary:
	
		\begin{compactitem}
			\item That may be seen distinctly; discoverable by the eye or the understanding; distinguishable.
		\end{compactitem}

	Why does Alma say this? He asks at the beginning of the verse, ``Isn't this real?'' I believe he's thinking of Korihor, who accused him and others of laboring under a ``foolish and a vain hope'' (\hyperref[ALMA30-13]{Alma 30:13}), having ``foolish traditions'' (\hyperref[ALMA30-14]{Alma 30:14}), of not being able to know of things that they ``do not see'' (\hyperref[ALMA30-15]{Alma 30:15}), and of being deluded about the reality of their experiences (\hyperref[ALMA30-16]{Alma 30:16}). So Alma then tells us about faith in \hyperref[ALMA-32-21]{verse 21} of this chapter: ``If ye have faith, ye hope for things which is not seen, which are true''---in other words, the things you are hoping for \textit{are real}, even though he admits that you can't see them. And now he gives an example involving this seed and the feelings and insights it creates inside you, and then he asks, isn't this real? How could it not be real? You're experiencing it right now! And then the discernible part. I read it not as saying that whatever is light is good \textit{because} it is discernible---I don't think he's saying, the explanation for the goodness is the discernibility. I think what's he saying is, you know it's \textit{real} because it's discernible---``whatsoever is light is good, [and you're able to make this judgment about the reality of the goodness] because it is discernible; therefore ye must know that it is good.'' The end part is really saying, ``To conclude that it really is good is a totally legitimate inference to make, based on what you have experienced.'' And then in the next verse, at the end of 36, he reminds us that we've really just been trying an experiment ``to know if the seed was good.'' And because our experiences, feelings, and insights are real---and we know they are real because they are discernible---then the experiment puts us in a position to make a good inference about the goodness of the seed.
	
	See also \hyperref[JOHN-3-19]{John 3:19--21}, 
	
	Some possibly relevant quotations:
	
		\begin{compactdesc}
			\item[Tertullian, ACCS NT 04a 130] The things that make us luminaries of the world are these---our good works. What is good, moreover (provided it is true and full), does not love darkness: it rejoices in being seen and exults over the very recognition it receives. To Christian modesty it is not enough to be so but to also appear so. For its fullness should be so great that it flows out from the mind to the clothing and bursts out from the conscience to the outward appearance.
			\item[Augustine, ACCS NT 04a 130] You must hate your own work and love the work of God in you. And when your own deeds begin to displease you, that is when your good works begin as you begin to find fault with your evil works. The beginning of good works is the confession of evil works, and then you do the truth and come to the light. How do you do the truth? You do not soothe or flatter yourself or say, ``I am righteous,'' while in actuality you are unrighteous. This is how you begin to do the truth. You come to the light so that your works may be shown to originate in God. And you have come to the light because this very sin in you, which displeases you, would not displease you if God did not shine on you and his truth show it to you. But the one who loves his sin, even after being admonished, hates the light admonishing him and flees from it so that his works that he loves may not be proved to be evil. 
		\end{compactdesc}
		
	% ----------------------
	\subsection{``Stretching forth'' one's hand in preaching and other activities}
	% ----------------------
		\label{EXSS-BOM-STRETCH-FORTH}
		
		In \hyperref[HELAMAN-13-4]{Helaman 13:4}, we read that Samuel gets on the wall and ``stretched forth his hand'' before speaking and prophesying. This study is about the different ways that ``stretching forth'' or ``stretching out'' someone's hand in this way is significant.
		
		First, a list of scriptures that mention ``stretching forth'' or ``stretching out'' someone's hand---if I'm counting correctly, there are 121 instances (with a search term of ``stretch* hand*'' in my scripture program):
		
		\begin{multicols}{3}
		
			\begin{compactitem}
				\item[{\hyperref[GENESIS-22-10]{Genesis 22:10}}]
				\item[{\hyperref[GENESIS-48-14]{Genesis 48:14}}]
				\item[{\hyperref[EXODUS-3-20]{Exodus 3:20}}]
				\item[{\hyperref[EXODUS-7-5]{Exodus 7:5}}]
				\item[{\hyperref[EXODUS-7-19]{Exodus 7:19}}]
				\item[{\hyperref[EXODUS-8-5]{Exodus 8:5}}]
				\item[{\hyperref[EXODUS-8-6]{Exodus 8:6}}]
				\item[{\hyperref[EXODUS-8-17]{Exodus 8:17}}]
				\item[{\hyperref[EXODUS-9-15]{Exodus 9:15}}]
				\item[{\hyperref[EXODUS-9-22]{Exodus 9:22}}]
				\item[{\hyperref[EXODUS-10-12]{Exodus 10:12}}]
				\item[{\hyperref[EXODUS-10-21]{Exodus 10:21}}]
				\item[{\hyperref[EXODUS-10-22]{Exodus 10:22}}]
				\item[{\hyperref[EXODUS-14-16]{Exodus 14:16}}]
				\item[{\hyperref[EXODUS-14-21]{Exodus 14:21}}]
				\item[{\hyperref[EXODUS-14-21]{Exodus 14:21}}]
				\item[{\hyperref[EXODUS-14-27]{Exodus 14:27}}]
				\item[{\hyperref[EXODUS-15-12]{Exodus 15:12}}]
				\item[{\hyperref[DEUTERONOMY-4-34]{Deuteronomy 4:34}}]
				\item[{\hyperref[DEUTERONOMY-5-15]{Deuteronomy 5:15}}]
				\item[{\hyperref[DEUTERONOMY-7-19]{Deuteronomy 7:19}}]
				\item[{\hyperref[DEUTERONOMY-11-12]{Deuteronomy 11:12}}]
				\item[{\hyperref[JOSHUA-8-18]{Joshua 8:18}}] twice
				\item[{\hyperref[JOSHUA-8-19]{Joshua 8:19}}]
				\item[{\hyperref[JOSHUA-8-26]{Joshua 8:26}}]
				\item[{\hyperref[1SAMUEL-24-6]{1 Samuel 24:6}}]
				\item[{\hyperref[1SAMUEL-26-9]{1 Samuel 26:9}}]
				\item[{\hyperref[1SAMUEL-26-11]{1 Samuel 26:11}}]
				\item[{\hyperref[1SAMUEL-26-23]{1 Samuel 26:23}}]
				\item[{\hyperref[2SAMUEL-1-14]{2 Samuel 1:14}}]
				\item[{\hyperref[2SAMUEL-24-16]{2 Samuel 24:16}}]
				\item[{\hyperref[1KINGS-8-42]{1 Kings 8:42}}]
				\item[{\hyperref[1CHRONICLES-21-16]{1 Chronicles 21:16}}]
				\item[{\hyperref[2CHRONICLES-6-32]{1 Chronicles 6:32}}]
				\item[{\hyperref[JOB-11-13]{Job 11:13}}]
				\item[{\hyperref[JOB-15-25]{Job 15:25}}]
				\item[{\hyperref[JOB-30-24]{Job 30:24}}]
				\item[{\hyperref[PSALMS-44-20]{Psalms 44:20}}]
				\item[{\hyperref[PSALMS-68-31]{Psalms 68:31}}]
				\item[{\hyperref[PSALMS-88-9]{Psalms 88:9}}]
				\item[{\hyperref[PSALMS-136-12]{Psalms 136:12}}]
				\item[{\hyperref[PSALMS-138-7]{Psalms 138:7}}]
				\item[{\hyperref[PSALMS-143-6]{Psalms 143:6}}]
				\item[{\hyperref[PROVERBS-1-24]{Proverbs 1:24}}]
				\item[{\hyperref[PROVERBS-31-20]{Proverbs 31:20}}]
				\item[{\hyperref[ISAIAH-5-25]{Isaiah 5:25}}] twice
				\item[{\hyperref[ISAIAH-9-12]{Isaiah 9:12}}] 
				\item[{\hyperref[ISAIAH-9-17]{Isaiah 9:17}}] 
				\item[{\hyperref[ISAIAH-9-21]{Isaiah 9:21}}] 
				\item[{\hyperref[ISAIAH-10-4]{Isaiah 10:4}}] 
				\item[{\hyperref[ISAIAH-14-26]{Isaiah 14:26}}] 
				\item[{\hyperref[ISAIAH-14-27]{Isaiah 14:27}}] 
				\item[{\hyperref[ISAIAH-23-11]{Isaiah 23:11}}] 
				\item[{\hyperref[ISAIAH-31-3]{Isaiah 31:3}}] 
				\item[{\hyperref[ISAIAH-45-12]{Isaiah 45:12}}] 
				\item[{\hyperref[JEREMIAH-6-12]{Jeremiah 6:12}}]
				\item[{\hyperref[JEREMIAH-15-6]{Jeremiah 15:6}}]
				\item[{\hyperref[JEREMIAH-32-21]{Jeremiah 32:21}}]
				\item[{\hyperref[JEREMIAH-51-25]{Jeremiah 51:25}}]
				\item[{\hyperref[EZEKIEL-6-14]{Ezekiel 6:14}}]
				\item[{\hyperref[EZEKIEL-10-7]{Ezekiel 10:7}}]
				\item[{\hyperref[EZEKIEL-14-9]{Ezekiel 14:9}}]
				\item[{\hyperref[EZEKIEL-14-13]{Ezekiel 14:13}}]
				\item[{\hyperref[EZEKIEL-16-27]{Ezekiel 16:27}}]
				\item[{\hyperref[EZEKIEL-20-33]{Ezekiel 20:33}}]
				\item[{\hyperref[EZEKIEL-20-34]{Ezekiel 20:34}}]
				\item[{\hyperref[EZEKIEL-25-7]{Ezekiel 25:7}}]
				\item[{\hyperref[EZEKIEL-25-13]{Ezekiel 25:13}}]
				\item[{\hyperref[EZEKIEL-25-16]{Ezekiel 25:16}}]
				\item[{\hyperref[EZEKIEL-30-25]{Ezekiel 30:25}}]
				\item[{\hyperref[EZEKIEL-35-3]{Ezekiel 35:3}}]
				\item[{\hyperref[DANIEL-11-42]{Daniel 11:42}}]
				\item[{\hyperref[HOSEA-7-5]{Hosea 7:5}}]
				\item[{\hyperref[ZEPHANIAH-1-4]{Zephaniah 1:4}}]
				\item[{\hyperref[ZEPHANIAH-2-13]{Zephaniah 2:13}}]
				\item[{\hyperref[MATTHEW-12-13]{Matthew 12:13}}] twice
				\item[{\hyperref[MATTHEW-12-49]{Matthew 12:49}}]
				\item[{\hyperref[MATTHEW-14-31]{Matthew 14:31}}]
				\item[{\hyperref[MATTHEW-26-51]{Matthew 26:51}}]
				\item[{\hyperref[MARK-3-5]{Mark 3:5}}] twice
				\item[{\hyperref[LUKE-6-10]{Luke 6:10}}]
				\item[{\hyperref[LUKE-22-53]{Luke 22:53}}]
				\item[{\hyperref[JOHN-21-18]{John 21:18}}]
				\item[{\hyperref[ACTS-4-30]{Acts 4:30}}]
				\item[{\hyperref[ACTS-12-1]{Acts 12:1}}]
				\item[{\hyperref[ACTS-26-1]{Acts 26:1}}]
				\item[{\hyperref[ROMANS-10-21]{Romans 10:21}}]
				\item[{\hyperref[1NEPHI-17-53]{1 Nephi 17:53}}] \textbf{1}
				\item[{\hyperref[1NEPHI-17-54]{1 Nephi 17:54}}] \textbf{1}
				\item[{\hyperref[2NEPHI-15-25]{2 Nephi 15:25}}] \textbf{2}
					\begin{compactitem}
						\item twice
						\item Compare Isaiah 5:25
					\end{compactitem}
				\item[{\hyperref[2NEPHI-19-12]{2 Nephi 19:12}}] \textbf{3}
					\begin{compactitem}
						\item Compare Isaiah 9:12 
					\end{compactitem} 
				\item[{\hyperref[2NEPHI-19-17]{2 Nephi 19:17}}] \textbf{3}
					\begin{compactitem}
						\item Compare Isaiah 9:17
					\end{compactitem}
				\item[{\hyperref[2NEPHI-19-21]{2 Nephi 19:21}}] \textbf{3}
					\begin{compactitem}
						\item Compare Isaiah 9:21
					\end{compactitem}
				\item[{\hyperref[2NEPHI-20-4]{2 Nephi 20:4}}] \textbf{3}
					\begin{compactitem}
						\item Compare Isaiah 10:4
					\end{compactitem}
				\item[{\hyperref[2NEPHI-24-26]{2 Nephi 24:26}}] \textbf{4}
					\begin{compactitem}
						\item Compare Isaiah 14:26
					\end{compactitem}
				\item[{\hyperref[2NEPHI-24-27]{2 Nephi 24:27}}] \textbf{4}
					\begin{compactitem}
						\item Compare Isaiah 14:27
					\end{compactitem} 
				\item[{\hyperref[JACOB-5-47]{Jacob 5:47}}] \textbf{5}
					\begin{compactitem}
						\item Note that this is master of the vineyard, who represents the Lord.
					\end{compactitem}
				\item[{\hyperref[JACOB-6-4]{Jacob 6:4}}] \textbf{3}
				\item[{\hyperref[MOSIAH-12-2]{Mosiah 12:2}}]
				\item[{\hyperref[MOSIAH-16-1]{Mosiah 16:1}}]
				\item[{\hyperref[ALMA-10-25]{Alma 10:25}}]
				\item[{\hyperref[ALMA-13-21]{Alma 13:21}}]
				\item[{\hyperref[ALMA-14-10]{Alma 14:10}}]
				\item[{\hyperref[ALMA-14-11]{Alma 14:11}}]
				\item[{\hyperref[ALMA-15-5]{Alma 15:5}}]
				\item[{\hyperref[ALMA-19-12]{Alma 19:12}}]
				\item[{\hyperref[ALMA-20-20]{Alma 20:20}}]
				\item[{\hyperref[ALMA-31-14]{Alma 31:14}}]
				\item[{\hyperref[ALMA-32-7]{Alma 32:7}}]
				\item[{\hyperref[HELAMAN-13-4]{Helaman 13:4}}]
				\item[{\hyperref[3NEPHI-11-9]{3 Nephi 11:9}}]
				\item[{\hyperref[3NEPHI-12-1]{3 Nephi 12:1}}]
				\item[{\hyperref[ETHER-3-6]{Ether 3:6}}]
				\item[{\hyperref[DC121-4]{DC 121:4}}]
				\item[{\hyperref[MOSES-7-36]{Moses 7:36}}]
				\item[{\hyperref[ABRAHAM-2-7]{Abraham 2:7}}]
				\item[{\hyperref[ABRAHAM-3-12]{Abraham 3:12}}]
			\end{compactitem}
		
		\end{multicols}
		
		\begin{compactenum}
				% 1
			\item To interact with other people, as commanded by God 
				% 2
			\item The Lord stretching forth his hand against people, negative connotation
				% 3
			\item The Lord's hand stretched out to people in mercy (often with ``still'' or for a long time)
				% 4
			\item The Lord stretching his hand over the nations or earth
				% 5
			\item A person stretching forth their hand to labor for salvation
		\end{compactenum}
		
		To perform an ordinance (sacrificial ordinance)
		To give a blessing
		To punish or take retribution (God acting)
		To execute a command from God (human acting)
		
		Think of significance of putting our arms up in the temple, or of the other arm positions; maybe we're invoking some of this